% problem_id: S001-lemma-factor-hom
% source_id: S001 (Stacks Project)
% source_file: local-cohomology.tex
% statement_locator: local-cohomology.tex lines 5113--5125
% proof_locator: local-cohomology.tex lines 5127--5231
% env: lemma
% id_note: label 在本来源内唯一
% pairing: tight (verified)
% status: complete (statement + author proof verbatim + reason + locators)
%
\begin{lemma}
\label{lemma-factor-hom}
With $q_0 = q(S)$ and $d = d(S)$ as above, let $M$ be an $(A, n, c)$-module
and let $\varphi : M \to I^n/I^{2n}$ be an $A$-linear map. Assume
$n \geq \max(q_0 + (1 + d)c, (2 + d)c)$ and if $d = 0$ assume
$n \geq q_0 + 2c$. Then the composition
$$
M \xrightarrow{\varphi} I^n/I^{2n} \to
I^{n - (1 + d)c}/I^{2n - (1 + d)c}
$$
is of the form $\sum a_i \psi_i$ with $a_i \in I^c$ and
$\psi_i : M \to I^{n - (2 + d)c}/I^{2n - (2 + d)c}$.
\end{lemma}

\begin{proof}
The case $d > 1$. Since we have a compatible system of maps
$p^*(I^q) \to \mathcal{O}_X(q)$ for $q \geq 0$ there are canonical maps
$p^*(I^q/I^{q + \nu}) \to \mathcal{O}_{Y_\nu}(q)$ for $\nu \geq 0$.
Using this and pulling back $\varphi$ we obtain a map
$$
\chi : p^*M \longrightarrow \mathcal{O}_{Y_n}(n)
$$
such that the composition
$M \to H^0(X, p^*M) \to H^0(X, \mathcal{O}_{Y_n}(n))$
is the given homomorphism $\varphi$ combined with the
map $I^n/I^{2n} \to H^0(X, \mathcal{O}_{Y_n}(n))$.
Since $\mathcal{O}_{Y_n}(n)$ is invertible on $Y_n$ the
linear map $\chi$ determines a section
$$
\sigma \in \Gamma(X, (p^*M)^\vee(n))
$$
with notation as in Remark \ref{remark-duals}.
The discussion in Remark \ref{remark-duals} shows
the cokernel and kernel of $can : p^*(M^\vee) \to (p^*M)^\vee$
are scheme theoretically supported on $Y_c$. By
Lemma \ref{lemma-annihilated} the map
$(p^*M)^\vee(n) \to (p^*M)^\vee(n - 2c)$ factors
through $p^*(M^\vee)(n - 2c)$; small detail omitted.
Hence the image of $\sigma$ in $\Gamma(X, (p^*M)^\vee(n - 2c))$
comes from an element
$$
\sigma' \in \Gamma(X, p^*(M^\vee)(n - 2c))
$$
By Lemma \ref{lemma-vanishing-coh-almost-projective} part (5),
the fact that $M^\vee$ is an $(A, n, c)$-module by
More on Algebra, Lemma \ref{more-algebra-lemma-dual-near-projective},
and the fact that $n \geq q_0 + (1 + d)c$ so $n - 2c \geq q_0 + (d - 1)c$
we see that the image of $\sigma'$ in $H^0(X, p^*M^\vee(n - (1 + d)c))$
is the image of an element $\tau$ in $I^{n - (1 + d)c}M^\vee$.
Write $\tau = \sum a_i \tau_i$ with $\tau_i \in I^{n - (2 + d)c}M^\vee$;
this makes sense as $n - (2 + d)c \geq 0$.
Then $\tau_i$ determines a homomorphism of modules
$\psi_i : M \to I^{n - (2 + d)c}/I^{2n - (2 + d)c}$
using the evaluation map $M \otimes M^\vee \to A/I^n$.

\medskip\noindent
Let us prove that this works\footnote{We hope some reader will suggest
a less dirty proof of this fact.}. Pick $z \in M$ and let us show that
$\varphi(z)$ and $\sum a_i \psi_i(z)$ have the same image in
$I^{n - (1 + d)c}/I^{2n - (1 + d)c}$.
First, the element $z$ determines a map
$p^*z : \mathcal{O}_{Y_n} \to p^*M$ whose composition with
$\chi$ is equal to the map $\mathcal{O}_{Y_n} \to \mathcal{O}_{Y_n}(n)$
corresponding to $\varphi(z)$ via the map
$I^n/I^{2n} \to \Gamma(\mathcal{O}_{Y_n}(n))$.
Next $z$ and $p^*z$ determine evaluation maps
$e_z : M^\vee \to A/I^n$ and $e_{p^*z} : (p^*M)^\vee \to \mathcal{O}_{Y_n}$.
Since $\chi(p^*z)$ is the section corresponding to $\varphi(z)$
we see that $e_{p^*z}(\sigma)$ is the section corresponding to $\varphi(z)$.
Here and below we abuse notation: for a map
$a : \mathcal{F} \to \mathcal{G}$ of modules on $X$
we also denote $a : \mathcal{F}(t) \to \mathcal{F}(t)$ the corresponding
map of twisted modules. The diagram
$$
\xymatrix{
p^*(M^\vee) \ar[d]_{can} \ar[r]_{p^*e_z} & \mathcal{O}_{Y_n} \ar@{=}[d] \\
(p^*M)^\vee \ar[r]^{e_{p^*z}} & \mathcal{O}_{Y_n}
}
$$
commutes by functoriality of the construction $can$. Hence
$(p^*e_z)(\sigma')$ in $\Gamma(Y_n, \mathcal{O}_{Y_n}(n - 2c))$
is the section corresponding to the image of $\varphi(z)$
in $I^{n - 2c}/I^{2n - 2c}$.
The next step is that $\sigma'$ maps to the image
of $\sum a_i \tau_i$ in $H^0(X, p^*M^\vee(n - (1 + d)c))$.
This implies that $(p^*e_z)(\sum a_i \tau_i) = \sum a_i p^*e_z(\tau_i)$
in $\Gamma(Y_n, \mathcal{O}_{Y_n}(n - (1 + d)c))$ is the section corresponding
to the image of $\varphi(z)$ in $I^{n - (1 + d)c}/I^{2n - (1 + d)c}$.
Recall that $\psi_i$ is defined from $\tau_i$ using an evaluation
map. Hence if we denote
$$
\chi_i : p^*M \longrightarrow \mathcal{O}_{Y_n}(n - (2 + d)c)
$$
the map we get from $\psi_i$, then we see by the same reasoning
as above that the section corresponding to $\psi_i(z)$ is
$\chi_i(p^*z) = e_{p^*z}(\chi_i) = p^*e_z(\tau_i)$. Hence we conclude that
the image of $\varphi(z)$ in $\Gamma(Y_n, \mathcal{O}_{Y_n}(n - (1 + d)c))$
is equal to the image of $\sum a_i\psi_i(z)$.
Since $n - (1 + d)c \geq q_0$ we have
$\Gamma(Y_n, \mathcal{O}_{Y_n}(n - (1 + d)c)) =
I^{n - (1 + d)c}/I^{2n - (1 + d)c}$ by Lemma \ref{lemma-bound-q-and-d} and
we conclude the desired compatibility is true.

\medskip\noindent
The case $d = 1$. Here we argue as above that we get
$$
\chi : p^*M \longrightarrow \mathcal{O}_{Y_n}(n),\quad
\sigma \in \Gamma(X, (p^*M)^\vee(n)),\quad
\sigma' \in \Gamma(X, p^*(M^\vee)(n - 2c)),
$$
and then we use
Lemma \ref{lemma-vanishing-coh-almost-projective} part (4)
to see that $\sigma'$ is the image of some element
$\tau \in I^{n - 2c}M^\vee$. The rest of the argument is the same.

\medskip\noindent
The case $d = 0$. Argument is exactly the same as in the
case $d = 1$.
\end{proof}
